\section{Discussion}
\label{sec:discussion}

The two-ratio description separates two questions that are often conflated in
hierarchical-memory optimisation: whether the active data can reside in the
fast tier, and whether compulsory transfer can be hidden behind computation.
The resulting labels are operational rather than universal phases.
$\theta_B=1$ follows exactly from the overlap definition, whereas
$\theta_C=0.5$ is a declared midpoint that makes ``majority resident'' precise.
A different application may choose another residency threshold without
changing Equation~\eqref{eq:lower_bound_regime}.

Both CPU probes support the qualitative capacity prediction. Their agreement
is useful because they stress different mechanisms: dependent pointer loads
probe the cache hierarchy directly, while the layered matrix harness combines
weight copying with computation. Neither test emulates accelerator DMA,
asynchronous streams, HBM, or collective communication. The data therefore do
not establish the I/O crossover on a GPU or the generality of the three labels
across model families.

The framework is most useful as a measurement discipline. Reporting $R_C$ and
$R_B$ alongside latency makes hidden assumptions about residency and overlap
inspectable. It also prevents a common circularity: using total step duration
in both the numerator of a bandwidth ratio and as the outcome being explained.
For deployment, $C_{\mathrm{fast}}$, $W$, $D$, sustained bandwidth, and isolated
compute time must be measured for the actual workload and hardware rather than
copied from nominal data-sheet values.

Several limitations are material. The present evidence comes from one CPU and
a small synthetic probe; it does not include production models, concurrent
requests, accelerator energy measurements, strategy comparisons, or unseen-
hardware classifier tests. The standard deviations describe repeated trials on
one machine and are not confidence intervals over a population of platforms.
The pointer-chain endpoints also span several hierarchy transitions, so their
ratio should not be interpreted as a universal effect size. Finally, the
additive decomposition in Equation~\eqref{eq:latency_decomp} is an accounting
model; overlap and contention couple its terms in real systems.

A decisive follow-up should preregister operating-point grids around
$R_C=0.5$ and $R_B=1$, collect raw event traces on at least two accelerator
architectures, and compare fixed orchestration policies without changing other
workload parameters. Such data could test boundary sharpness, policy ranking,
and generalisation. Until those measurements are available, the contribution
of this work is the corrected analytical formulation, open protocol, and CPU
proof of concept rather than a claim of universal accelerator validation.
